\documentclass[10pt]{article}
\usepackage[margin=20mm]{geometry}
\usepackage{amsmath,amssymb,mathtools,fontspec,unicode-math,xurl,hyperref}
\IfFontExistsTF{DejaVu Serif}{\setmainfont{DejaVu Serif}}{\setmainfont{Latin Modern Roman}}
\setmathfont{latinmodern-math.otf}
\setlength{\parindent}{0pt}\setlength{\parskip}{6pt}\sloppy
\title{Technical review of
Bitangents of real algebraic curves:
signed count and constructions}
\author{AI-assisted technical review}
\date{}
\begin{document}
\maketitle
Abstract: The reviewed source is arXiv:2402.03993v1 of 6 February 2024. The signed Euler-characteristic calculation and the algebraic counts in the real deformation argument survive the checks recorded below. The asserted extension to every field of characteristic different from two does not: a rational, admissible two-section gives a rational quartic family whose unique bitangents near two specified rational limiting chords are not rational at arbitrarily small rational parameters. Smoothness, ordinary contact, local uniqueness, and one generic specialization are certified independently. This is Major for the stated arithmetic extension, without refuting the real main results. A separate impossible root-count prescription in Lemma 4.7 is Minor; a realizable replacement is supplied. The retained source evidence images have now been inspected. Scientific adjudication remains pending.\par
\section*{Identity, version, and scope}
The source is T. Blomme, E. Brugallé and C. Garay, Bitangents of real algebraic curves: signed count and constructions, arXiv:2402.03993v1, submitted 6 February 2024 at 13:45:19 UTC [source]. The associated journal DOI is `10.4171/JEMS/1716`. The publisher lists acceptance on 11 April 2025 and online publication on 7 November 2025 [publisher]. The acquired 38-page PDF is the arXiv version, not a verified copy of the accepted or journal text. The frozen identifier and year remain M25-0339 and 2025. Every page was read as extracted text, including the references. Page references below are one-based PDF pages, agreeing with the printed pagination.\par
The audit follows the hypotheses of genericity, the sign convention, the Hilbert-square stratification, the real deformation count, the stated construction counts, and the arithmetic assertion in Section 5. It includes independent exact symbolic calculations and analytic proofs. The bounded continuation visually inspected all 35 retained source evidence pages, including Figures 1–22. The displayed schemes, labels and cited proof interfaces were checked; this does not independently reprove every Brusotti/Harnack realization. The external existence and regularity theorems used by the authors have not all been reproved.\par
A bounded search of the arXiv record, publisher record, author-related publication records, and title-based correction/code/supplement queries did not locate a separate paper-specific implementation, supplement, or correction notice. This is not proof of absence. No author code was run, and absence of code is not a defect. The publisher requires a subscription for its article text; no equivalence to that unacquired text is asserted.\par
\section*{Definitions, proof chain, and checks without a finding}
A generic plane curve here is smooth and has dual singularities consisting only of ordinary nodes and cusps (pp. 1–2). For degree \(d=2k\), write \(p,n\) for the numbers of even and odd real ovals. The bitangent sign is unchanged if the two complementary sign regions are exchanged. Real bitangents with conjugate contact points contribute positively; for two real contacts the sign is the product of the two local side signs. These distinctions matter in passing from an unsigned count to \(t_s\).\par
\subsection*{Euler-characteristic bookkeeping}
For the double plane \(S\) branched along the curve, put \(a=p-n\). The identities used on pp. 9–14 are \[\chi(\mathbb CS)=6+d(d-3),\qquad
 \chi(\mathbb RS)=2a,\qquad
 \chi(\mathbb RP^+)=a,\quad \chi(\mathbb RP^-)=1-a.\] Compactly supported Euler characteristic is used on strata. The real sign regions are two-dimensional manifolds, so their ordinary and compactly supported Euler characteristics agree. There is no sign discrepancy from that passage.\par
For a nodal curve of arithmetic genus \(g\), with \(k_+\) hyperbolic real nodes, \(k_-\) elliptic real nodes, and \(k_{\mathbb C}\) conjugate pairs of non-real nodes, set \(r=k_--k_+\). The normalization and local Hilbert-square contributions give \[\chi(\mathbb RC^{[2]})
 =\frac{r^2+(2-2g+k_++k_-+2k_{\mathbb C})}{2}-k_+-k_-
 =1-g+k_{\mathbb C}-k_++\binom{k_--k_+}{2}.\] The cusp term is \(1-g\). These formulas reproduce the two integrations in the paper. For the projection from pairs on \(S\), the distinct-real-pair part contributes \[\frac{\chi(\mathbb RS)^2-\chi(\mathbb RS)}2-\chi(\mathbb RP^+)
 =-2a(1-a),\] and the conjugate-pair part contributes \[\frac{\chi(\mathbb CS)-\chi(\mathbb RS)}2-\chi(\mathbb RP^-)
 =2+k(2k-3).\] The tangent-direction strata have zero Euler contribution. Comparison with integration over lines yields \[t_s=2(p-n)(p-n-1)+\frac{d(d-2)}2.\] The integer product \((p-n)(p-n-1)\) is nonnegative, including the cases of empty real locus and negative \(p-n\). The complex and real analogues are consistent with the stated Plücker and Klein counts. These algebraic identities were independently reproduced exactly, not inferred from numerical samples.\par
\subsection*{Degeneration and counting}
The \(k\)-section has Newton triangle with vertices \((0,0)\), \((0,k)\) and \((4k,0)\), is smooth in the affine chart, and has nonrepeated factors on the upper edge (Definition 3.3, pp. 17–18). The projection general-position condition on p. 22 requires distinct simple critical values, exactly \[N=4k(k-1)=d(d-2)\] in number. It does not prohibit concurrency among unrelated limiting chords.\par
Proposition 3.6 constructs the degree-\(2k\) family by replacing the normal coordinate with \((y^2-xz)/t\). The local dual-curve argument in Theorem 3.9 gives one nearby bitangent for each chord joining distinct critical values, and \(2(k-2)\) near each tangent to the conic. The algebraic total is \[\binom{N}{2}+2(k-2)N
 =\frac{d(d-2)(d-3)(d+3)}2,\] the full complex bitangent count. The distinct-chord branches are simple intersections in the dual degeneration.\par
For the coalescing case, the local model can be checked by writing a ramified branch as \(x=u^4+t\alpha(u)\), \(y=u^2+\cdots\) with \(\alpha'(0)\ne0\), and an unramified branch with a different value \(\widehat\alpha(0)\). After \(t=s^2\), \(u=s\xi\), \(v=s\eta\), the leading tangent equations are \[\eta=\frac{\alpha'(0)}{4\xi},\qquad
 \eta^2=\widehat\alpha(0)-\alpha(0).\] The roots are simple when the two heights differ. The sign condition after reversing the normal-coordinate sign is precisely the real-height rule in Theorem 3.9. This verifies the local mechanism and its hypotheses; the global geometric existence arguments retain their stated dependencies.\par
The arithmetic of the sextic count on p. 30 gives \[276+2\{12-(p+n-1)\}+4(p-1)=298+2(p-n).\] The Harnack-family sum on p. 31 is \[\sum_{i=1}^{k-2}4i(k-1-i)=\frac{2k(k-1)(k-2)}3;\] doubling gives the number of non-real bitangents. Both were checked independently. The special sextic with 318 real bitangents depends additionally on the displayed realizations; the displayed topology has now been inspected, while its full algebraic realization remains dependent on the cited construction inputs.\par
\section*{Arithmetic specialization does not preserve the ground field}
Problem 1 (A rational limiting chord can have no nearby rational bitangent).  Section 5(7), p. 34, asserts that Theorem 3.9 has a weak version over every field \(K\) of characteristic different from two: a limiting chord defined over \(K\) deforms to a bitangent defined over \(K\) of \(C_t\) for \(t\in K\). The following example over \(K=\mathbb Q\) violates that specialization assertion, including arbitrarily small rational parameters.\par
\subsection*{An admissible two-section and the prescribed family}
Take \[\mathcal A=\left\{\frac12,2,\frac13,3\right\},\qquad
 P(u)=\prod_{a\in\mathcal A}(u^2-a^2),\qquad C:\ w^2=P(u).\] All eight roots \(\pm a\) are distinct rational numbers. At a root \(r\), the projection has the expansion \(u=r+w^2/P'(r)+O(w^4)\), so its eight critical points are nondegenerate with distinct critical values. The affine curve is smooth. The constant term is nonzero and the upper-edge equation is \(w^2-u^8\), with distinct factors. In the chart at infinity, \(s=1/u\), \(W=w/u^4\), the equation is \[W^2=\prod_{a\in\mathcal A}(1-a^2s^2).\] At \(s=0\) the two points \(W=\pm1\) are smooth and unramified. Thus the compactified two-section also meets the required boundary conditions.\par
Equation reference: quartic. Proposition 3.6 gives exactly the rational quartic family \[F_t(x,y,z)=(y^2-xz)^2-t^2\prod_{a\in\mathcal A}(x-a^2z)=0.
 \] On the limiting conic, use \((x,y,z)=(u^2,u,1)\). Two of its specified rational chords are \[L_1:\ y=\frac25(x+z),\quad (u=\tfrac12,2),
 \qquad
 L_2:\ y=\frac3{10}(x+z),\quad (u=\tfrac13,3).\] The reciprocal symmetry does not violate any stated projection assumption.\par
\subsection*{Exact bitangent equations}
Equation reference: restriction, mbranches. Set \(x=U+V\), \(z=U-V\), and \(c=625/9\). Direct multiplication gives \[\prod_{a\in\mathcal A}(x-a^2z)
 =16U^4-cU^2V^2+cV^4.\] Restrict [quartic] (\#quartic) to \(y=2mU\), and put \(A=4m^2-1\). The restriction is \[(1-ct^2)V^4+(2A+ct^2)U^2V^2+(A^2-16t^2)U^4.
 \] For \(0<|t|<3/25\) this is a quadratic square exactly when \[cA^2+cA+16+\frac{c^2-64c}{4}t^2=0.\] Its discriminant in \(A\) is \((175/9)^2(1-ct^2)\). Consequently the two positive-slope branches satisfy \[m_\pm(t)^2=\frac18\pm\frac7{200}
                 \sqrt{1-\frac{625}{9}t^2},
 \qquad
 m_+(0)=\frac25,\quad m_-(0)=\frac3{10}.
 \] These are exact algebraic equations, independently reproduced in the retained calculation.\par
For every integer \(N\ge105\), \[(3N-1)^2<9N^2-625<(3N)^2.\] Hence \(\sqrt{1-625/(9N^2)}\) is irrational. Formula [mbranches] (\#mbranches) shows that \(m_\pm(1/N)^2\), and therefore \(m_\pm(1/N)\), are irrational. A line \(y=m(x+z)\) is defined over \(\mathbb Q\) exactly when its normalized coefficient \(m\) is rational. Both displayed bitangents fail that condition.\par
\subsection*{No second, nonsymmetric rational branch is available}
This uniqueness point is essential. A general nearby line has the form \(y=hU+\ell V\). Set \[A=h^2-1,\qquad B=2h\ell,\qquad D=1+\ell^2.\] Its restriction is \[(AU^2+BUV+DV^2)^2-t^2(16U^4-cU^2V^2+cV^4).\] If it is a bitangent, this binary quartic is a square \((aU^2+bUV+eV^2)^2\) over the algebraic closure. Comparison of coefficients gives \[ab=AB,\quad be=BD,\quad
 a^2=A^2-16t^2,\quad e^2=D^2-ct^2.\] In neighborhoods of \(L_1,L_2\), the numbers \(h,A,D\) are nonzero. If \(\ell\ne0\), division gives \(a/e=A/D\), hence, for \(t\ne0\), \[cA^2=16D^2.\] At the two limiting chords, respectively, \(cA^2-16D^2\) equals \(-7\) and \(112/9\). It remains nonzero on fixed neighborhoods. Thus every nearby bitangent has \(\ell=0\). The quadratic equation above has one separated root near each corresponding value \(A=-9/25,-16/25\), and the positive limiting \(h\) fixes its square root. Each specified chord has exactly the one nearby branch already computed. This argument also holds in complex neighborhoods.\par
\subsection*{Smoothness and two ordinary real contacts}
The quartic is a quadratic form in \((U^2,V^2,y^2)\), with matrix \[Q_t=\begin{pmatrix}
 1-16t^2&-1+ct^2/2&-1\\
 -1+ct^2/2&1-ct^2&1\\
 -1&1&1
 \end{pmatrix}.\] For \(0<t<3/25\), its first-order principal minors are nonzero. Its second-order principal minors are \[-16t^2-\frac{30625}{324}t^4,\qquad -16t^2,\qquad -ct^2,\] and its determinant is \(-30625t^4/324\). If a projective singular point existed, its nonzero-coordinate support would give a nonzero vector in the kernel of a principal submatrix of \(Q_t\). The listed minors exclude this. The quartic is therefore smooth over \(\mathbb C\) throughout this interval.\par
For \(t=1/N\), \(N\ge105\), both roots satisfy \[-\frac{16}{25}<A<-\frac9{25},\quad
 1-ct^2>0,\quad 2A+ct^2<0.\] Equation [restriction] (\#restriction) becomes \[(1-ct^2)\left(V^2+
       \frac{2A+ct^2}{2(1-ct^2)}\,U^2\right)^2.\] The quadratic has two distinct real roots. Since the curve is smooth, these give two distinct contacts of intersection multiplicity exactly two. The lines are actual ordinary real bitangents, not formal candidates or hyperflexes.\par
\subsection*{An independently certified generic specialization}
Take \(t=1/1000\). The non-square integer above is \(8999375\), between \(2999^2=8994001\) and \(3000^2=9000000\). The seven principal minors at this parameter are \[\frac{62499}{62500},\quad \frac{14399}{14400},\quad 1,\quad
 -\frac{8294449}{518400000000},\quad
 -\frac1{62500},\quad-\frac1{14400},\quad
 -\frac{49}{518400000000}.\] To verify the stronger generic-dual hypothesis, denominators in \(F_{1/1000}\) were cleared and its homogeneous Hessian determinant \(H\) computed independently. A good reduction modulo 13 gives \[\overline F=-5U^4-4U^2V^2-5U^2y^2-5V^4+5V^2y^2-4y^4.\] On \(y=1\), the Gröbner basis of \[\left(\overline F,\overline H,\,
 \partial_U\overline F\,\partial_V\overline H-
 \partial_V\overline F\,\partial_U\overline H\right)\] is \((1)\). On \(y=0,U=1\), the ideal consisting of \(\overline F,\overline H\) and all three projective gradient minors also has basis \((1)\). At the only omitted projective point \([0:1:0]\), \(\overline F=8\ne0\). These exact finite-field computations certify transverse intersection of the quartic and its Hessian. The projective bad-intersection locus is closed over the integers localized away from the denominators; properness implies that a characteristic-zero bad point would specialize to a bad point here. There is none.\par
Thus the characteristic-zero quartic has no hyperflex. A smooth plane quartic without hyperflex has dual with ordinary nodes and cusps: a line cannot be tangent at three distinct points by Bézout, nor be both a flex tangent and an additional tangent, and distinct contact points give distinct tangent directions on the dual. The specialization meets the source’s genericity requirement. This is a deterministic certificate, not a floating-point discriminant test.\par
The same one-parameter algebraic family has a proper closed set of nongeneric parameters, since a generic member has just been certified. That set is finite. Infinitely many of the rational parameters \(1/N\) therefore remain generic and approach zero, while the irrationality and uniqueness arguments apply to all \(N\ge105\). Restricting to sufficiently small generic parameters cannot repair the assertion.\par
\subsection*{Impact and repair boundary}
The conclusion fails for the precise rational family [quartic] (\#quartic) and rational limiting chords satisfying the stated assumptions. No claim is made that this curve has no other rational bitangents. The counterexample concerns the asserted lifting of each specified chord.\par
There is a meaningful weaker statement: the two branches [mbranches] (\#mbranches) belong to \(\mathbb Q[[t]]\), since the inner square root has constant term \(1\) and the outer roots have nonzero rational constants \(2/5,3/10\). Formal or appropriate Henselian lifting does not imply rationality at a rational specialization. An analytic real branch also gives no such implication. Replacing the assertion by a formal/Henselian statement, with its hypotheses, or restricting it to suitable analytic fields changes its arithmetic content.\par
This is Major for the claimed arbitrary-field extension and its proposed use to bound \(K\)-bitangents. The real deformation and counting theorems are not refuted by it. The location in Perspectives does not make the false arithmetic conclusion a notation error. No correction in the unacquired journal text, or author adoption of this repair boundary, is asserted.\par
\section*{The non-real root prescription}
Problem 2 (The stated root count can be odd).  In the proof of Lemma 4.7, p. 33, the real degree-\(k\) perturbing polynomials \(R_1,R_2\) are required to have simple roots and \(\lfloor k/2\rfloor\) non-real roots. For \(k=2\) or \(3\), this requires exactly one non-real root. A real polynomial has conjugate non-real roots with equal multiplicities, so the requirement is impossible. The same obstruction occurs for \(k\equiv2,3\pmod4\).\par
The natural correction is \(\lfloor k/2\rfloor\) pairs of non-real roots. More is needed than that wording to justify the subsequent assertion of zero or two real critical points of \(R_1/R_2\). The following explicit choices supply that property.\par
Choose distinct small positive numbers \(a_j,b_j\) with \(a_jb_j<1\), avoiding common roots, and write \[H(z)=\prod_{j=1}^{m}\frac{z+a_j}{1+b_jz}.\] Every factor is positive for \(z\ge0\) and has derivative \((1-a_jb_j)/(1+b_jz)^2>0\). Thus \(H'(z)>0\) for \(m>0\).\par
For \(k=2m+1\), take \(R_1(\theta)=\theta\prod_j(\theta^2+a_j)\) and \(R_2(\theta)=\prod_j(1+b_j\theta^2)\) as homogeneous binary forms of degree \(k\). The second form has its single real root at infinity. Then \[\frac{d}{d\theta}\left(\theta H(\theta^2)\right)
 =H(\theta^2)+2\theta^2H'(\theta^2)>0.\] At infinity the map is asymptotic to a nonzero constant times \(\theta\), so it is unramified there as well. A sufficiently small real Möbius change of the domain moves the pole at infinity to a finite point and gives affine polynomials of exact degree \(k\), still arbitrarily close to \(\theta^k\) and \(1\). It preserves simple roots, their real/non-real status, and absence of real critical points. Each polynomial has \(2m\) non-real roots and one real root. The case \(k=1\) is the identity map with the same conclusion.\par
For \(k=2m\), omit the leading factor \(\theta\) and use the same products. All \(k\) roots of each polynomial are non-real and simple. The quotient is \(H(\theta^2)\), whose only real critical points on \(\mathbb RP^1\) are \(0\) and infinity. The derivative at zero vanishes to first order because \(H'(0)>0\). At infinity, \[H(\theta^2)=H_\infty
 \left(1-\Bigl[\sum_j(1/b_j-a_j)\Bigr]\theta^{-2}
       +O(\theta^{-4})\right),\] and the coefficient is nonzero, so the critical point there is also simple. All these choices can approach the original pair \(\theta^k,1\).\par
This verifies a local algebraic replacement for the root prescription and the asserted real critical-point counts. Persistence and chosen smoothing of all the rational curve’s nodes additionally use the source’s geometric construction and Brusotti input; the displayed diagrams have now been inspected, while the external existence and smoothing results remain dependencies. The finding is Minor: the impossible prescription can be replaced under the original degree and reality conditions, and no counterexample to Lemma 4.7 or Proposition 1.6 is established.\par
\section*{Further clarifications and limits}
The unnumbered proposition on pp. 31–32 states \(r_k=k^2+O(k)\), whereas its proof specifies \(r_k=4k(k-1)\) for odd \(k\), and at least \(4(k-1)(k-2)\) for even \(k\). Together with the total ramification bound \(4k(k-1)\), these displayed proof counts give \(r_k=4k^2+O(k)\). This stronger estimate yields \(\binom{r_k}{2}=8k^4+O(k^3)\), matching the leading complex bitangent count in degree \(2k\). The printed weaker leading coefficient alone would give only \(k^4/2+O(k^3)\) and would not imply Theorem 1.12. The normalization should be made consistent. This is recorded as a text-level clarification, not as a proven counterexample to an existential proposition or to the complete connectedness construction.\par
For Proposition 5.1, the leading conic count is compatible with Gathmann’s primary formula: the enumerative polynomial has leading term \(d^{10}/5!\) [gathmann]. The source’s \(9k(k-1)\) branch values for degree \(d=3k\) give \(\binom{9k(k-1)}5=d^{10}/5!+O(d^9)\). This checks the leading arithmetic only. The full real cubic-normal-bundle construction and all conic transversality claims in the Perspectives section have not been independently certified. The unrelated questions posed there are not treated as established theorems or as defects.\par
The \(k\)-section and local analytic conditions have been kept distinct from the stronger claim that every arithmetic specialization preserves a ground field. No arbitrary positive-characteristic theorem, no full independent graphical-realization proof, and no journal-version result is certified by this report.\par
\section*{Conclusion}
The overall assessment is Major, confined to the false arithmetic specialization assertion in Section 5(7), with the separate Minor issue in Lemma 4.7. The signed-count identity and checked real local-deformation algebra are not contradicted. The real graphical construction claims retain the explicit verification limits above. Scientific adjudication remains `pending`; no human review or independent acceptance is recorded.\par
9 T. Blomme, E. Brugallé and C. Garay, Bitangents of real algebraic curves: signed count and constructions, arXiv:2402.03993v1, 6 February 2024. https://arxiv.org/abs/2402.03993v1. EMS Press, article metadata for DOI 10.4171/JEMS/1716, consulted 27 October 2026. https://ems.press/journals/jems/articles/14299255. A. Gathmann, The number of plane conics that are five-fold tangent to a given curve, Compositio Math. 141 (2005), 487–501, DOI 10.1112/S0010437X04001083. The leading enumerative formula on p. 488 was checked; the full relative Gromov–Witten proof was not rederived.\par
\section*{Supporting checks}
46ecd9a262e787f66d7d6b78cf098bd65243da592b1c175b5c7cf67681cb63e5.\par
Paper identifier: `M25-0339`.\par
The original initial executor read all 38 source pages through actual Poppler extraction. Poppler returned exit code zero with an xref-reconstruction warning, which is retained; the source bytes were unchanged. The executor wrote and ran the independent SymPy calculations: the Euler/counting identities, the quartic square condition, the non-square certificate, all Ciani principal minors, the finite-field Hessian certificate, the general-line uniqueness obstruction, and the derivative identity for the local repair. Both actual calculation runs returned zero and their source, stdout, stderr, commands, and receipts are retained. No author code was executed.\par
The global finite-field genericity calculation and the explicit monotonicity repair were performed by the initial executor. Final reasoning and grading are the initial executor’s responsibility.\par
No personally observed UI state or backend equivalence is invented. Main supplied scheduling and byte handling only.\par
The report’s text and image checks are separate: a parsed or rendered page is not a viewed image.\par
The native bounded continuation inspected all 35 retained source evidence pages, confirmed the printed hypotheses and both existing findings, and retained the original exact calculations without rerunning them. Its exact-source cached final build and all report-page inspections are recorded separately under this paper’s recovery evidence. The Gathmann consultation remains an attributed original execution fact; no external primary page is credited as newly viewed. Scientific decision remains pending, with no human review or independent acceptance recorded.\par
\end{document}
